\section{Título de la sección}

% ── Caso de estudio ───────────────────────────────────────────────────────────
\begin{caso}{Nombre del caso (año)}
Enunciado del caso tal como lo entrega el profesor.
\end{caso}

\begin{fuentes}
  \fuente{https://ejemplo.org/articulo}
\end{fuentes}
% \nocite{clave-del-bib}   % añade la entrada en tex/bib/refs.bib y cítala aquí

% ── Preguntas abiertas ────────────────────────────────────────────────────────
\pregunta{Texto de la pregunta.}
\begin{respuesta}
\end{respuesta}

% ── Selección múltiple ────────────────────────────────────────────────────────
% \opcion*{...} marca la respuesta elegida (admite varias).
\begin{mcq}{Enunciado de la pregunta de selección múltiple (1 o más respuestas válidas):}
  \opcion{Primera opción.}
  \opcion{Segunda opción.}
  \opcion{Tercera opción.}
\end{mcq}

% ── Laboratorio con evidencias ────────────────────────────────────────────────
% 1) Comando o payload en un listing (estilo global "cysec", ya aplicado):
\begin{lstlisting}[language=bash]
comando --con-sus-flags "objetivo"
\end{lstlisting}
% 2) Captura: si existe img/ejemplo.png se incluye; si no, caja punteada
%    con el nombre esperado y warning «Captura pendiente» en el .log.
\captura{ejemplo.png}{Título descriptivo de la captura}
% 3) El análisis va después de la captura; vacío ⇒ caja naranja y warning
%    «Análisis pendiente».
\begin{analisis}
\end{analisis}
